\documentclass{article}
\usepackage{tcee-repaso}
\tcetema[Economía del desarrollo]{3.B.28}{Economía de los países en desarrollo. Teorías recientes del desarrollo económico. Evidencia empírica, con especial referencia a la aproximación experimental, e implicaciones para el diseño de políticas.}
\tcefecha{11/11/2026}
\begin{document}
\cabecerarepaso
\begin{columnas}

% ==========================================================================
\bloque{Esquema del cante}
\esquemacante{%
\apartadocante{Introducción: «unos nacen con estrella\ldots»}{2,5}
\apartadocante{I. Hechos estilizados}{3}
\subcante{1. Medición del desarrollo}{1,5}
\subcante{2. Características de los países en desarrollo}{1,5}
\apartadocante{II. Aproximación macroeconómica}{10}
\subcante{1. Tendencias: crecimiento, desigualdad, pobreza}{1,5}
\subcante{2. Teorías pioneras y su replanteamiento}{5,5}
\subcante{3. Teorías recientes: complementariedades, instituciones}{3}
\apartadocante{III. Aproximación microeconómica}{13,5}
\subcante{1. Pioneros: \autor{Sen} y \autor{Deaton}}{2}
\subcante{2. Comportamiento y mercados de los pobres}{3}
\subcante{3. La aproximación experimental}{6}
\subcante{4. Implicaciones para las políticas}{2,5}
\apartadocante{Conclusión}{1}}
\emph{Gráficos del cante}: \autor{Lewis} y la trampa de pobreza (\hyperref[rep:II]{II}), en 3 pasos cada uno; función clave: el sesgo de selección (\hyperref[rep:III]{III}). Pegatinas: «una golondrina no hace verano» (gran empujón), «café para todos» (Consenso), «el hábito no hace al monje» (2.º óptimo), «ver para creer» (RCT), «del dicho al hecho\ldots» (validez externa), «obras son amores\ldots» (políticas).

% ==========================================================================
\bloque[2,5]{Introducción}
\begin{itemize}
\item \textbf{Enganche}: «Unos nacen con estrella y otros estrellados»: la estrella es el país. Kenia, esperanza de vida 64 años; España, 84; renta pc PPA ×8. \concepto{Prima de ciudadanía} (\autor{Milanovic}, 2016). ¿Por qué unos siguen pobres y qué sabemos de cómo salir?
\item \textbf{Relevancia}: \emph{pobreza} $<$3\,\$/día: 43,4\,\% (1990) $\to$ 10,1\,\% (2024), 824 M; 7 de cada 10 en África subsahariana. \emph{Convergencia parcial}: PIB pc emergentes ×2,6 (1995-2024) vs. ×1,5 avanzadas (FMI). \emph{Nobel}: \autor{Myrdal} (1974), \autor{Lewis}-\autor{Schultz} (1979), \autor{North} (1993), \autor{Sen} (1998), \autor{Deaton} (2015), \autor{Banerjee}-\autor{Duflo}-\autor{Kremer} (2019, experimental), AJR (2024, instituciones): los 2 extremos, programa concreto y reglas del juego.
\item \textbf{Contextualización}: \concepto{crecimiento} = ↑ sostenido renta pc (cuantitativo) \vertema{3.A.43}; \concepto{desarrollo} = transformación productiva, social e institucional que amplía las \emph{capacidades} (\autor{Sen}, 1999; cualitativo). Crecimiento: necesario, no suficiente (Guinea Ecuatorial y Sri Lanka, $\approx$12.700\,\$ PPA; IDH puesto 133 vs. 89). \concepto{País en desarrollo}: renta baja o media (BM), $>$80\,\% de la población. Historia: \autor{Smith} (1776), \autor{Solow} (1956); economía del desarrollo tras la 2.ª GM y la descolonización; 3 etapas: optimismo intervencionista (50-60), contrarrevolución neoclásica (80), desde los 90 formalización + revolución experimental. Sin programa asentado: estructural (teoría $\to$ datos) vs. experimental (datos $\to$ teoría) \vertema{3.A.1}.
\item \textbf{Problemática}: ¿qué comparten y cómo se mide el atraso? ¿Por qué una teoría distinta de la del crecimiento y qué queda? ¿Qué explican las teorías recientes y qué dice la evidencia experimental? ¿Qué implica para las políticas?
\item \textbf{Estructura}: \hyperref[rep:I]{I}, hechos estilizados; \hyperref[rep:II]{II}, aproximación macro (tendencias; por qué algunos países se quedan atrás: de los pioneros a las instituciones); \hyperref[rep:III]{III}, aproximación micro (\autor{Sen} y \autor{Deaton}, comportamiento y mercados, experimentos, políticas).
\end{itemize}

% ==========================================================================
\bloque[3]{I. Hechos estilizados\label{rep:I}}
\begin{ideaclave}
Países muy distintos con rasgos medibles que se refuerzan entre sí. Ninguna cifra resume el desarrollo: el IDH es la medida de referencia, la renta su base material y la de la clasificación, y la pobreza su cara más dura.
\end{ideaclave}

\sub{I.1. Medición del desarrollo\label{rep:I.1}}
\begin{itemize}
\item \concepto{IDH} (PNUD, 1990; capacidades de \autor{Sen}) \vertema{3.A.24}: salud (esperanza de vida), educación (años esperados y medios), nivel de vida ($\ln$ RNB pc PPA). Normalizadas 0-1 y \emph{media geométrica} (desde 2010): una dimensión baja no se compensa del todo.
\begin{formula}
\centering $\mathrm{IDH}=(I_S\cdot I_E\cdot I_R)^{1/3},\quad I_R=\dfrac{\ln(\mathrm{RNB}/L)-\ln100}{\ln75.000-\ln100}$
\end{formula}
\item[] Niveles: muy alto $\geq$0,8; alto 0,7; medio 0,55; bajo $<$0,55 (casi todos en África subsahariana). 2023: Islandia 0,972; España 28.ª (0,918); Sudán del Sur 0,388.
\item \concepto{RNB pc} (PIB + rentas netas del exterior: las filiales repatrían beneficios): base del IDH y de la \emph{clasificación BM} (1 de julio, método Atlas, desde jul. 2026): baja $\leq$1.175\,\$ (25); media-baja 1.176-4.635 (47); media-alta 4.636-14.375 (59); alta $>$14.375 (86). En \concepto{PPA}: ricos/pobres ×65 a tipos de mercado, ×26 en PPA. No basta: es un medio y una media.
\item \concepto{Pobreza absoluta} (mínimo fijo) vs. \concepto{relativa} (UE: $<$60\,\% de la mediana). \concepto{Pobreza extrema}: $<$3,00\,\$ PPA 2021 desde jun. 2025 (antes 2,15 de 2017); con 4,20 y 8,30\,\$, 18\,\% y 46\,\%.
\item \concepto{Índice de pobreza multidimensional} (PNUD-Oxford, 2010, sustituye a los de pobreza humana): 1.100 M, la mitad menores.
\end{itemize}

\sub{I.2. Características de los países en desarrollo\label{rep:I.2}}
\begin{itemize}
\item \textbf{Económicos}: \emph{dualidad}: agricultura 59\,\% del empleo en renta baja, 3\,\% en alta; \emph{subempleo} e informalidad (58\,\% del empleo mundial, OIT) más que paro abierto; \emph{poco ahorro y finanzas poco profundas}: crédito privado 13\,\% vs. 152\,\% del PIB; \emph{dependencia de primarios} (ley de \autor{Engel}: elasticidad renta $<$1); \concepto{enfermedad holandesa} (\autor{Corden}-\autor{Neary}, 1982): el recurso aprecia el TCR y expulsa industria y agricultura; \emph{vulnerabilidad} (deuda, clima) \vertema{3.B.29} \vertema{3.B.31}; \emph{desigualdad} alta (América Latina, África austral).
\item \textbf{Institucionales}: Estado con poca capacidad (impuestos $\approx$10\,\% del PIB en renta baja; aranceles, señoreaje); reglas débiles (propiedad, corrupción, inestabilidad).
\item \textbf{Demográficos}: \emph{transición a medias}: renta baja, esperanza de vida 65 años, 4,5 hijos, población +2,7\,\%/año, 41\,\% $<$15 años $\Rightarrow$ $n$ alta ↓ renta de EE (\autor{Solow}) \vertema{3.A.43}.
\item Los rasgos se refuerzan: de síntomas a causas (macro) y a restricciones de los hogares (micro).
\end{itemize}

% ==========================================================================
\bloque[10]{II. Aproximación macroeconómica\label{rep:II}}
\begin{ideaclave}
A escala de los países, el atraso no se corrige solo. Los pioneros acertaron en las preguntas (cambio estructural, equilibrios múltiples, fallos de coordinación) sin apenas formalizarlas; sus recetas, intervencionistas y luego liberales, fallaron por universales. Las teorías recientes las formalizan y apuntan a las instituciones.
\end{ideaclave}

\sub{II.1. Tendencias: crecimiento, desigualdad y pobreza\label{rep:II.1}}
\begin{itemize}
\item \emph{Crecimiento}: sin convergencia en el s.\,XX; desde 2000, algo y frágil (\autor{Kremer}-\autor{Willis}-\autor{You}, 2022) \vertema{3.A.45}; \concepto{trampa de la renta media} (\autor{Gill}-\autor{Kharas}, 2007). IDH: todas las regiones mejoran, África subsahariana la última.
\item \emph{Evolución}: 43,4 $\to$ 10,1\,\% (1990-2024), lenta desde 2014. Asia oriental 67 $\to$ 2\,\%, meridional 50 $\to$ 3\,\%; China, casi 2/3 de la caída (940 M, 83\,\%, $\to$ $\approx$0 desde 2019). África subsahariana 62 $\to$ 45\,\%, pero 321 $\to$ 581 M: 70\,\% de los pobres.
\end{itemize}
\graficorepaso[\linewidth]{pobreza-regiones}{Pobreza extrema por regiones, 1990-2026 (millones; tasa mundial, eje dcho.): Asia se vacía, África crece}
\begin{itemize}
\item \concepto{Curva de \autor{Kuznets}} (1955): U invertida renta-Gini; la dualidad de \autor{Lewis}. \emph{No es ley}: 150 países, $R^2=0{,}10$, la U la empujan América Latina y África austral; dentro de cada país, rara vez (\autor{Deininger}-\autor{Squire}, 1998); \autor{Piketty} (2014): caída de 1914-45 por guerras, inflación e impuestos \vertema{3.A.45}.
\item \concepto{Curva del elefante} (\autor{Lakner}-\autor{Milanovic}, 2016): 1988-2008 ganan el centro mundial y el 1\,\% más rico; pierden las clases medias-bajas ricas. Desigualdad \emph{entre} países ↓ desde los 90, \emph{dentro} ↑.
\end{itemize}

\sub{II.2. Teorías pioneras y su replanteamiento\label{rep:II.2}}
\begin{itemize}
\item \autor{Solow} (1956) \vertema{3.A.43}: EE con progreso exógeno y \concepto{convergencia condicional} \vertema{3.A.45}. Insatisfacción: (1) explica poco: residuo 87,5\,\% (EE.\,UU. 1909-49; \autor{Solow}, 1957), «medida de nuestra ignorancia» (\autor{Abramovitz}, 1956); (2) mundo nivelador, pero los pobres no se acercaban.
\item \autor{Lewis}, \autor{Rosenstein-Rodan}, \autor{Nurkse}, \autor{Myrdal}, \autor{Kaldor}, \autor{Hirschman}, \autor{Prebisch}, \autor{Singer}: «heterodoxos» de inspiración keynesiana. (1) \emph{Convergencia no garantizada}: círculos viciosos y equilibrios múltiples; (2) \emph{cambio estructural como causa}: economía dual, la industria es el motor.
\item Excepción: \autor{Schultz} (1964): campesinos pobres pero eficientes; faltan factores nuevos. Germen de la contrarrevolución.
\item Dualidad $\to$ \emph{etapa} (\autor{Lewis}, \autor{Rostow}); no convergencia $\to$ \emph{situación} (\autor{Myrdal}, \autor{Kaldor}, \autor{Nurkse}, CEPAL).
\item \textbf{\autor{Lewis} (1954)}, \concepto{modelo dual}: ahorrar del 4-5 al 12-15\,\% sube la parte de los capitalistas. Tradicional: PMg$_L\approx0$, salario institucional $\overline w_A$ (producto medio); moderno: $w_M=\overline w_A(1+\text{prima})$, \concepto{oferta ilimitada de trabajo} (perfectamente elástica). $L_M$ donde $p\,\mathrm{PMg}_L(K)=w_M$; el beneficio se reinvierte, $D(K)$ $\to$ dcha., más empleo \emph{al mismo salario}; \concepto{punto de \autor{Lewis}}: se agota el excedente, $w$ ↑ en ambos sectores.
\end{itemize}
\graficorepaso[0.78\linewidth]{lewis}{\autor{Lewis}: oferta horizontal en $w_M$; la reinversión desplaza $D(K)$ hasta el punto de \autor{Lewis}}
\begin{itemize}
\item[] Implicaciones: motor = acumulación del beneficio (oferta); desigualdad ↑ al principio (\autor{Kuznets}); no subir salarios antes de tiempo. \emph{Extensiones}: \autor{Ranis}-\autor{Fei} (1961): el excedente se agota antes, precio de alimentos ↑ $\Rightarrow$ crecimiento equilibrado con el agro; \autor{Harris}-\autor{Todaro} (1970), paro urbano:
\begin{formula}
\centering $w_A=\dfrac{L_M}{L_U}\,\overline w_M$\quad (salario urbano \emph{esperado})
\end{formula}
\item[] crear empleo urbano puede ↑ el paro. \emph{Evidencia}: Japón, Corea, China (debate del punto de \autor{Lewis} hacia 2010; \autor{Zhang} et al., 2011); éxodo rural español de los 60 \vertema{4.A.27}. \emph{Críticas}: reinversión incompleta o fuga, técnica ahorradora de trabajo, movilidad perfecta.
\item \textbf{\autor{Rostow} (1960)}, 5 etapas: sociedad tradicional, condiciones previas, \emph{despegue} (inversión $<$5 $\to$ $>$10\,\% de la renta, sectores líderes), madurez, alto consumo de masas. \concepto{Brecha de financiación} (\autor{Domar}) \vertema{3.B.29}:
\begin{formula}
\centering $\dfrac{I}{Y}=v\,\bar\gamma_Y\ \Rightarrow\ \text{ayuda «necesaria»}=v\,\bar\gamma_Y-s$
\end{formula}
\item[] Críticas: descripción (historia británica); \autor{Gerschenkron} (1962): los atrasados se saltan etapas con más Estado o bancos; la ayuda no fue a inversión ni esta a crecimiento; \autor{Domar} renegó (\autor{Easterly}, 1999).
\item \textbf{\autor{Myrdal} (1957)}, \concepto{causación circular acumulativa} (precedente: \autor{Young}, 1928): \emph{efectos de retroceso} (la región rica atrae capital y trabajo) vs. \emph{de difusión}; si dominan los primeros, divergencia $\Rightarrow$ política regional \vertema{3.B.6}. Sin modelo: no dice cuál domina.
\item \textbf{\autor{Kaldor} (1966, 1970)}, \concepto{ley de \autor{Kaldor}-\autor{Verdoorn}}: productividad ↑ con la producción (escala dinámica, aprendizaje). \autor{Dixon}-\autor{Thirlwall} (1975):
\begin{formula}
\centering $\gamma_{Y/L}=a_V+b_V\gamma_Y,\ \ \gamma_Y=a_D+b_D\gamma_{Y/L}\ \Rightarrow\ \gamma_Y^*=\dfrac{a_D+b_Da_V}{1-b_Db_V}$
\end{formula}
\item[] $b_Db_V<1$, estable; $>1$, inestable: espiral virtuosa o viciosa. Modelo de demanda, sin acumulación.
\item \textbf{\autor{Nurkse} (1953)}, \concepto{círculos viciosos}: oferta (renta baja $\to$ poco ahorro $\to$ poco $K$ $\to$ baja productividad) y demanda (mercado pequeño $\to$ poca inversión) $\Rightarrow$ \concepto{trampa de pobreza}. \autor{Solow} con producción en S (\autor{Azariadis}-\autor{Drazen}, 1990): $sf(\tilde k)/\tilde k$ corta 3 veces $n+g+\delta$: $\tilde k_B$ (trampa) y $\tilde k_A$ estables, $\tilde k_U$ umbral inestable. \concepto{Gran empujón} más allá de $\tilde k_U$, luego se retira. Variante: \concepto{trampa demográfica} (\autor{Nelson}, 1956).
\end{itemize}
\graficorepaso[0.78\linewidth]{trampa-pobreza}{Trampa de pobreza: $\tilde k_B$ y $\tilde k_A$ estables, $\tilde k_U$ umbral; el gran empujón lleva más allá de $\tilde k_U$}
\begin{itemize}
\item \autor{Rosenstein-Rodan} (1943): fábrica de zapatos sola no rentable; muchas a la vez, sí: \concepto{fallo de coordinación}, «una golondrina no hace verano»; crecimiento \emph{equilibrado}. \autor{Murphy}-\autor{Shleifer}-\autor{Vishny} (1989): técnica moderna con coste fijo $F$, $c<1$ y prima $w>1$; sola vende $q^A=L/N$, todas $q^B=(L/N-F)/c$:
\begin{formula}
\centering $B(q)=(1-c\,w)\,q-w\,F;\quad \text{2 equilibrios si } B(q^A)<0\leq B(q^B)$
\end{formula}
\end{itemize}
\graficorepaso[0.7\linewidth]{gran-empujon}{MSV: con $w_2$ la empresa sola pierde ($A$) y todas juntas ganan ($B$); $w_1$, se moderniza sola; $w_3$, nunca}
\begin{itemize}
\item[] Sin prima ($w=1$), un solo equilibrio. Política: subvención temporal, plan de inversión o expectativas coordinadas. \autor{Hirschman} (1958): crecimiento \emph{desequilibrado}, sectores con más \concepto{eslabonamientos} hacia atrás y hacia delante.
\item[] \emph{Evidencia}: entre países, pocas trampas (\autor{Kraay}-\autor{McKenzie}, 2014); en hogares, más plausibles. Justificó grandes programas que fracasaron.
\item \textbf{CEPAL} (\autor{Prebisch}, 1949; \autor{Singer}, 1950): centro-periferia, «correa de transmisión rota»; \concepto{RRI} $=p_X/p_M$ ↓ a LP (\autor{Engel}, proteccionismo agrario; sindicatos y poder de mercado en el centro) $\Rightarrow$ \concepto{ISI} (industria naciente \vertema{3.B.8}; racionar divisas). \concepto{Ley de \autor{Thirlwall}} (1979): $\gamma_Y=(\varepsilon_{X,Y^*}/\varepsilon_{M,Y})\gamma_{Y^*}$; \emph{dependencia} (\autor{Frank}, 1967; \autor{Cardoso}-\autor{Faletto}, 1969). Balance: RRI ↓ débil y con rupturas; ISI cara, búsqueda de rentas (\autor{Krueger}, 1974), India 6,5\,\% $\to$ 5\,\% de la renta pc de EE.\,UU. (1950-80); a favor, \autor{Chang} (2002). Hoy vuelve \vertema{3.B.9}.
\item \textbf{Orientación a la exportación}: 4 tigres; manufacturas intensivas en trabajo, TC realista, insumos libres, apoyos condicionados a exportar; 6-7\,\% anual; también ahorro, educación, estabilidad, reforma agraria.
\item \emph{60-70}: reconsideración conservadora (\autor{Bauer}, H.\,\autor{Johnson}): libre comercio, agricultura, fallos del gobierno; ayuda perjudica (\autor{Bauer}, 1972). \emph{70}: \concepto{necesidades básicas} (OIT, Ginebra, 1976). \emph{80}: \concepto{contrarrevolución neoclásica} (\autor{Krueger}, \autor{Bhagwati}, \autor{Balassa}): deuda 1982, agotamiento ISI, tigres; ajuste FMI-BM.
\item \concepto{Consenso de Washington} (\autor{Williamson}, nov. 1989), 10 puntos en 4 bloques: estabilidad macro (disciplina fiscal, gasto, reforma tributaria); precios de mercado (tipos de interés, TC competitivo); apertura (comercio, IDE); mercado y Estado (privatización, desregulación, propiedad).
\item \emph{Crisis}: México 94, Asia 97, Rusia 98, Argentina 2001. «\emph{Café para todos}», secuencia; lectura parcial de Asia (BM, 1993; \autor{Stiglitz}, 1996); 2.º óptimo \vertema{3.A.22}; equidad. Correcciones: \emph{market-friendly} (BM, 1991), instituciones (\autor{Burki}-\autor{Perry}, 1998), \concepto{Consenso de Santiago}. Queda la macro: necesaria, no suficiente.
\end{itemize}

\sub{II.3. Teorías recientes: complementariedades e instituciones\label{rep:II.3}}
\begin{itemize}
\item \textbf{Crecimiento endógeno}: $AK$ (\autor{Rebelo}, 1991) \vertema{3.A.44}, sin convergencia garantizada; competencia imperfecta $\Rightarrow$ espacio para la política:
\begin{formula}
\centering $\gamma_{Y/L}=s\,A-(n+\delta)$
\end{formula}
\item[] \concepto{Paradoja de \autor{Lucas}} (1990): el capital no va a los pobres ($h$, PTF, instituciones, riesgo) \vertema{3.B.29}.
\item \textbf{Fallos de coordinación}: el gran empujón formalizado (MSV, 1989): «nadie invierte porque nadie invierte»; de intuición a resultado.
\item \textbf{\concepto{Teoría del anillo O}} (\autor{Kremer}, 1993; \emph{Challenger}, 1986): $m$ tareas, éxito $q_i$:
\funcionclave{\text{valor esperado}\propto\overline V\prod_{i=1}^{m}q_i}
\item[] 10 tareas: 0,95 $\to$ 60\,\% del valor; 0,90 $\to$ 35\,\%. Emparejamiento por cualificación, salarios convexos (fuga de cerebros), productos sencillos, no formarse es racional; infraestructuras malas = anillo defectuoso.
\item \textbf{\concepto{Distancia a la frontera}} (\autor{Acemoglu}-\autor{Aghion}-\autor{Zilibotti}, 2006): lejos, \emph{imitar} (grandes grupos, bancos); cerca, \emph{innovar} (competencia, entrada, mercados de capitales). \concepto{Trampa de la renta media}: desde 1990, 34 a renta alta ($>$1/3 por UE o petróleo); 108 siguen, 6.000 M (BM, 2024): invertir, \emph{incorporar}, \emph{innovar} (Corea).
\item \concepto{Instituciones}: «reglas del juego» (\autor{North}, 1990): fijan incentivos y quién se queda los frutos. \autor{Hall}-\autor{Jones} (1999): infraestructura social. \concepto{Inclusivas} vs. \concepto{extractivas} (\autor{Acemoglu}-\autor{Robinson}, 2012).
\item AJR (2001): mortalidad de colonos como instrumento (explica el \emph{nivel}); VI 0,94 $>$ MCO 0,52 \vertema{3.A.45}. Persisten por \emph{conflicto social} (AJR, 2005). Botsuana (1966): límites a los jefes, propiedad, diamantes sin enfermedad holandesa (AJR, 2003).
\item ¿Cuáles? Sin consenso: China (empresas de municipios, doble vía de precios). \concepto{Instituciones de segundo óptimo} (\autor{Rodrik}, 2008) \vertema{3.A.22}: «el hábito no hace al monje», importa la función.
\item \concepto{Diagnóstico del crecimiento} (\autor{Hausmann}-\autor{Rodrik}-\autor{Velasco}, 2008): la restricción con mayor precio sombra. Inversión baja por: bajo rendimiento social, baja \emph{apropiabilidad} o \emph{financiación} cara (señal: tipos reales altos). \concepto{Autodescubrimiento} (\autor{Hausmann}-\autor{Rodrik}, 2003): premiar el descubrir y retirar si no hay resultados.
\end{itemize}

% ==========================================================================
\bloque[13,5]{III. Aproximación microeconómica\label{rep:III}}
\begin{ideaclave}
A escala de las personas, la pobreza cambia cómo se decide y deja a los hogares fuera de mercados que fallan. La evidencia experimental ha cambiado el diseño de las políticas, sin sustituir a la teoría ni a la macro.
\end{ideaclave}

\sub{III.1. Pioneros: \autor{Sen} y \autor{Deaton}\label{rep:III.1}}
\begin{itemize}
\item Precedente: \autor{Schultz} (1964), campesinos pobres pero eficientes.
\item \textbf{\autor{Sen}} (Nobel 1998): \concepto{capacidades}, pobreza = privación (base del IDH y del IPM); \concepto{titularidades}: Bengala 1943, hambruna sin caída de la producción (\autor{Sen}, 1981): acceso, no disponibilidad; las democracias con prensa libre no tienen grandes hambrunas.
\item \textbf{\autor{Deaton}} (Nobel 2015): \emph{encuestas de hogares} para medir consumo, pobreza y efectos de precios (\autor{Deaton}, 1997), base de las cifras del BM; calorías en la India ↓ con crecimiento (\autor{Deaton}-\autor{Drèze}, 2009). Después, crítico de los RCT.
\end{itemize}

\sub{III.2. Comportamiento y mercados de los pobres\label{rep:III.2}}
\begin{itemize}
\item \concepto{Brecha de aspiraciones} (\autor{Ray}, 2006): ni pequeña ni enorme; solo la intermedia mueve a invertir; reproduce la desigualdad (\autor{Genicot}-\autor{Ray}, 2017). \emph{Desesperanza} (\autor{Duflo}, 2012).
\item \emph{Escasez}: cultivadores de caña indios razonan peor antes de la cosecha (\autor{Mani} et al., 2013): «la pobreza agota» (\autor{Mullainathan}-\autor{Shafir}, 2013).
\item \emph{Trampas de los hogares} (renta en S: nutrición, \autor{Dasgupta}-\autor{Ray}, 1986; activos; crédito). BRAC, Bangladés: sobre el umbral de activos, salen de la pobreza 11 años (\autor{Balboni} et al., 2022); entre países, raras. Política: contar con el comportamiento \vertema{3.A.1}; con umbral, ayuda grande y temporal.
\item \textbf{Mercados que fallan}. \emph{Crédito}: sin colateral, información asimétrica $\Rightarrow$ racionamiento (\autor{Stiglitz}-\autor{Weiss}, 1981) \vertema{3.A.13}. \concepto{Responsabilidad conjunta} (Grameen, \autor{Yunus}, Nobel de la Paz 2006): selección (\autor{Ghatak}, 1999), vigilancia (\autor{Stiglitz}, 1990), sanciones sociales. \emph{Seguros}: redes informales cubren choques idiosincrásicos, no agregados (sequía). M-Pesa (Kenia): 194.000 hogares, 2\,\% (\autor{Suri}-\autor{Jack}, 2016; cuasiexperimental).
\end{itemize}

\sub{III.3. La aproximación experimental\label{rep:III.3}}
\begin{itemize}
\item \textbf{Método}: \concepto{RCT} (\autor{Kremer}, Kenia; \autor{Banerjee}-\autor{Duflo}, India): «ver para creer». \emph{Resultados potenciales} (\autor{Rubin}, 1974): $Y_i(1)$, $Y_i(0)$, solo se ve uno; \concepto{ATE} $\equiv\mathbb{E}[Y_i(1)-Y_i(0)]$. \concepto{Sesgo de selección}:
\funcionclave{\begin{aligned}&\underbrace{\mathbb{E}[Y_i|D_i{=}1]-\mathbb{E}[Y_i|D_i{=}0]}_{\text{observado}}=\underbrace{\mathbb{E}[Y_i(1)-Y_i(0)|D_i{=}1]}_{\text{efecto en tratados}}\\&\qquad+\underbrace{\mathbb{E}[Y_i(0)|D_i{=}1]-\mathbb{E}[Y_i(0)|D_i{=}0]}_{\text{sesgo de selección}}\end{aligned}}
\item[] Libros de texto: padres implicados $\Rightarrow$ exagera. \emph{Al azar}: sesgo nulo, diferencia de medias = ATE (\autor{Angrist}-\autor{Pischke}, 2009); en MCO, $Y_i=b_0+b_1D_i+u_i$.
\item[] \concepto{Validez interna} (sin sesgo en la muestra) vs. \concepto{externa} (fuera; la muestra no es aleatoria: \autor{Allcott}, 2015, 111 evaluaciones de ahorro energético, las primeras sobrestiman). \concepto{Intención de tratar} si no todos cumplen; contagios; escala (\autor{Duflo}-\autor{Glennerster}-\autor{Kremer}, 2007). \emph{Cuasiexperimentos}: dif. en dif. (\autor{Duflo}, 2001, escuelas en Indonesia), regresión discontinua, VI, experimentos naturales \vertema{4.B.6}.
\item \textbf{Evidencia} (J-PAL, 2003: $>$2.700 evaluaciones en $>$140 países). \emph{Salud}: desparasitación en Busia, absentismo $-1/4$, también en vecinos (\autor{Miguel}-\autor{Kremer}, 2004), 20 años después más consumo y salario; tasa pequeña $\Rightarrow$ participación $-80$\,\% (\autor{Kremer}-\autor{Miguel}, 2007); mosquiteras (\autor{Cohen}-\autor{Dupas}, 2010); vacunas en la India 6 $\to$ 18 $\to$ 39\,\% con lentejas, a mitad de coste. \emph{Educación}: libros solo a los buenos; comidas +31\,\% asistencia; funciona enseñar según el nivel (Pratham) e incentivos (cámaras: absentismo ÷2). \emph{Crédito}: 6 países, efectos modestos y poca demanda (\autor{Banerjee}-\autor{Karlan}-\autor{Zinman}, 2015); sin responsabilidad conjunta, no más impagos (\autor{Giné}-\autor{Karlan}, 2014). \emph{Agricultura}: fertilizante y \emph{sesgo al presente}: descuento tras la cosecha (\autor{Duflo}-\autor{Kremer}-\autor{Robinson}, 2011). \emph{Transferencias}: condicionadas (\emph{Progresa}, 1997), incondicionales (GiveDirectly; multiplicador $\approx$2,5, poca inflación), \emph{graduación} (6 países). Nobel 2019.
\item \textbf{Valoración}. \emph{Ventajas}: credibilidad, política basada en evidencia (coste por resultado), diseño antes que evaluación. \emph{Críticas}: validez externa sin teoría (\autor{Deaton}, 2010; \autor{Deaton}-\autor{Cartwright}, 2018) («del dicho al hecho\ldots»); mecanismos y reparto (\autor{Heckman}-\autor{Smith}, 1995), modelos estructurales; jerarquía de la evidencia (\autor{Ravallion}, 2020); pensar en pequeño, «arreglar una ventana en una casa que se cae»; \autor{Rodrik} (2009): diagnóstico + experimento; Hawthorne (\autor{Levitt}-\autor{List}, 2011); ética. Balance: no explican el crecimiento, sí qué programas funcionan.
\end{itemize}

\sub{III.4. Implicaciones para las políticas\label{rep:III.4}}
\begin{enumerate}
\item \emph{Primero, el diagnóstico} (HRV, 2008): ni gran empujón ni Consenso para todos.
\item \emph{Macro e instituciones, condición necesaria}, adaptadas al país y a su distancia a la frontera.
\item \emph{Evaluar antes de escalar}: ONG $\neq$ administración (\autor{Banerjee} et al., 2017); J-PAL, IPA, BM, cooperación española \vertema{4.A.26} \vertema{4.B.6}.
\item \emph{Diseñar para personas reales}: precio cero en prevención, incentivos pequeños e inmediatos, compromisos, trámites sencillos.
\item \emph{Dar dinero} y, con umbrales, un empujón grande (graduación); ayuda ligada a evidencia e instituciones, no a la brecha \vertema{3.B.29} \vertema{3.B.30} \vertema{3.B.42}.
\item \emph{Política industrial moderna}: fallos concretos, apoyos condicionados y retirada (\autor{Juhász}-\autor{Lane}-\autor{Rodrik}, 2024); invertir $\to$ incorporar $\to$ innovar.
\end{enumerate}
\begin{itemize}
\item Marco: Agenda 2030, ODS 1; clima, responsabilidades comunes pero diferenciadas \vertema{3.B.31}. «Obras son amores, que no buenas razones»: pensar en grande para el marco, en pequeño para el diseño.
\end{itemize}

% ==========================================================================
\bloque[1]{Conclusión}
\begin{itemize}
\item \textbf{Recapitulación}: «En conclusión,» rasgos comunes y 3 medidas (IDH, renta, pobreza); macro: pobreza 43 $\to$ 10\,\%, concentrada en África; los pioneros acertaron en 3 ideas (muchos factores, se refuerzan, margen para la política) pero sus recetas generales fallaron; formalización e instituciones como causa última; micro: \autor{Sen}, \autor{Deaton} y experimentos, en debate con lo estructural.
\item \textbf{Relevancia}: 1 de cada 10, $\approx$800 M en pobreza extrema; 108 países en renta media.
\item \textbf{Extensiones}: \vertema{3.A.43} \vertema{3.A.44} \vertema{3.A.45} \vertema{3.B.9} \vertema{3.B.29} \vertema{3.B.30} \vertema{3.B.42} \vertema{4.A.26} \vertema{3.B.31} \vertema{4.B.6} \vertema{3.A.1}.
\item \textbf{Opinión}: los 2 Nobel son complementarios: China y Corea crecieron con ahorro (China 51\,\% del PIB en 2008), inversión, exportaciones y Estado coordinador, no con experimentos; dentro de ese marco, saber cómo responden las personas multiplica cada euro.
\item \textbf{Idea final}: «unos nacen con estrella y otros estrellados»: Kenia o España sigue siendo estrella, pero no destino (Botsuana, China, Corea).
\end{itemize}

% ==========================================================================
\bloque{Autores y fechas}
\hito{1943}{\autor{Rosenstein-Rodan}: gran empujón}
\hito{1949}{\autor{Prebisch}; 1950, \autor{Singer}: RRI ↓}
\hito{1953}{\autor{Nurkse}: círculos viciosos}
\hito{1954}{\autor{Lewis}: oferta ilimitada de trabajo}
\hito{1955}{\autor{Kuznets}: U invertida}
\hito{1956}{\autor{Solow}; \autor{Nelson}: trampa demográfica}
\hito{1957}{\autor{Myrdal}: causación acumulativa}
\hito{1958}{\autor{Hirschman}: crecimiento desequilibrado}
\hito{1960}{\autor{Rostow}: etapas}
\hito{1961}{\autor{Ranis}-\autor{Fei}}
\hito{1962}{\autor{Gerschenkron}: los atrasados se saltan etapas}
\hito{1964}{\autor{Schultz}: campesinos eficientes}
\hito{1966}{\autor{Kaldor}: ley de \autor{Verdoorn}}
\hito{1970}{\autor{Harris}-\autor{Todaro}: paro urbano}
\hito{1974}{\autor{Krueger}: búsqueda de rentas; \autor{Rubin}: resultados potenciales}
\hito{1976}{OIT: necesidades básicas}
\hito{1979}{\autor{Thirlwall}: restricción externa}
\hito{1989}{\autor{Williamson}: Consenso; MSV: gran empujón formal}
\hito{1990}{\autor{North}: reglas del juego; PNUD: IDH; \autor{Lucas}: paradoja}
\hito{1993}{\autor{Kremer}: anillo O; BM: milagro asiático}
\hito{2001}{AJR: mortalidad de colonos; \autor{Duflo}: Indonesia}
\hito{2004}{\autor{Miguel}-\autor{Kremer}: desparasitación}
\hito{2006}{\autor{Acemoglu}-\autor{Aghion}-\autor{Zilibotti}; \autor{Ray}: aspiraciones}
\hito{2007}{\autor{Gill}-\autor{Kharas}: renta media}
\hito{2008}{HRV: diagnóstico; \autor{Rodrik}: 2.º óptimo}
\hito{2010}{\autor{Deaton}: crítica de los RCT}
\hito{2015}{\autor{Banerjee}-\autor{Karlan}-\autor{Zinman}: microcréditos}
\hito{2019}{Nobel a \autor{Banerjee}, \autor{Duflo}, \autor{Kremer}}
\hito{2022}{\autor{Balboni} et al.: umbral de activos}
\hito{2024}{Nobel a AJR; BM: trampa de la renta media}

% ==========================================================================
\bloque{Datos clave}
\begin{itemize}
\item Pobreza extrema ($<$3\,\$ PPA 2021): \dato{43,4\,\% (1990) $\to$ 10,1\,\% (2024), 824 M}{BM, PIP 2026}; África subsahariana \dato{70\,\% de los pobres, 581 M}{BM}.
\item PIB pc 1995-2024: \dato{×2,6 emergentes vs. ×1,5 avanzadas}{FMI, 2026}.
\item Renta pc alta/baja: \dato{×65 a TC de mercado, ×26 en PPA}{BM, WDI}.
\item Umbral renta baja \dato{1.175\,\$}{BM, jul. 2026}; países: \dato{25 baja, 47 media-baja, 59 media-alta, 86 alta}{BM}.
\item IDH 2023: España \dato{0,918 (28.ª)}{PNUD, 2025}; Islandia \dato{0,972}{PNUD}; Sudán del Sur \dato{0,388}{PNUD}.
\item Pobreza multidimensional: \dato{1.100 M}{PNUD-Oxford, 2025}.
\item Empleo agrario: \dato{59\,\% renta baja, 3\,\% alta}{BM, 2025}; empleo informal \dato{58\,\%}{OIT, 2023}.
\item Crédito privado: \dato{13\,\% vs. 152\,\% del PIB}{BM}; renta baja: \dato{4,5 hijos, +2,7\,\%/año, 41\,\% $<$15}{BM}.
\item Renta media: \dato{34 llegan a alta desde 1990; 108 siguen, 6.000 M}{BM, 2024}.
\item J-PAL: \dato{$>$2.700 evaluaciones, $>$140 países}{J-PAL, 2026}.
\end{itemize}

% ==========================================================================
\bloque{Preguntas del tribunal}
\begin{enumerate}
\item \textbf{¿Crecimiento o desarrollo?} Cuantitativo vs. cualitativo (capacidades de \autor{Sen}); necesario, no suficiente (Guinea Ecuatorial y Sri Lanka).
\item \textbf{¿Por qué media geométrica en el IDH?} Una dimensión baja no se compensa del todo; desde 2010.
\item \textbf{¿Se cumple \autor{Kuznets}?} No como ley: $R^2=0{,}10$, empujada por América Latina y África austral; rara dentro de cada país.
\item \textbf{¿Por qué \autor{Harris}-\autor{Todaro} predice paro urbano?} Se emigra hasta igualar el salario \emph{esperado}; crear empleo urbano atrae más emigrantes.
\item \textbf{¿Existen las trampas de pobreza?} Entre países, apenas (\autor{Kraay}-\autor{McKenzie}); en hogares sí, con umbral de activos (BRAC).
\item \textbf{¿Por qué aleatorizar?} Elimina el sesgo de selección: tratados y controles iguales en promedio; diferencia de medias = ATE.
\item \textbf{Principal crítica a los RCT}: validez externa (\autor{Deaton}); además mecanismos, ética, pensar en pequeño. Respuesta: replicar en varios países y ligar a teoría.
\item \textbf{¿Qué quedó del Consenso de Washington?} La macro (disciplina fiscal, inflación baja, TC realista): necesaria, no suficiente.
\item \textbf{¿Por qué el capital no va a los pobres?} Paradoja de \autor{Lucas} (1990): diferencias de $h$, productividad, instituciones y riesgo.
\item \textbf{¿Qué es el diagnóstico del crecimiento?} Buscar la restricción con mayor precio sombra (rendimiento social, apropiabilidad, financiación) y reformar ahí (HRV, 2008).
\item \textbf{¿Instituciones: cuáles?} Sin receta: segundo óptimo (\autor{Rodrik}), función y no forma; diagnóstico del crecimiento.
\end{enumerate}

\begin{trampa}
\begin{itemize}
\item Pobreza extrema: 3,00\,\$ PPA \emph{2021} (antes 2,15 de 2017); clasificación BM a tipos de mercado (Atlas), no en PPA.
\item IDH: media \emph{geométrica} (aritmética antes de 2010); renta en $\ln$; educación con años esperados y medios.
\item \autor{Lewis}: oferta \emph{perfectamente elástica} a $w_M>\overline w_A$; el motor es el ahorro de los capitalistas (oferta), no la demanda.
\item \autor{Harris}-\autor{Todaro}: el paro está en la ciudad; \autor{Ranis}-\autor{Fei}: el excedente se agota antes.
\item Trampa de pobreza: $\tilde k_B$ y $\tilde k_A$ estables, $\tilde k_U$ \emph{inestable}; no la predice \autor{Solow}: exige producción en S.
\item \autor{Rosenstein-Rodan}/\autor{Nurkse}: crecimiento \emph{equilibrado}; \autor{Hirschman}: \emph{desequilibrado} (eslabonamientos).
\item MSV: sin prima salarial ($w=1$) no hay equilibrios múltiples.
\item \autor{Kaldor}: $b_Db_V>1$ inestable (acumulativo); $<1$ estable.
\item Anillo O: cualificaciones que se \emph{multiplican}: 0,95$^{10}\approx$0,60 vs. 0,90$^{10}\approx$0,35.
\item AJR: VI 0,94 $>$ MCO 0,52; el instrumento explica el \emph{nivel}.
\item Aleatorizar da validez \emph{interna}, no externa.
\item Tasa pequeña por antiparasitarios: participación $-80$\,\%; cobrar no hace valorar más.
\item M-Pesa: estimación cuasiexperimental, no RCT. Microcréditos: efectos modestos.
\end{itemize}
\end{trampa}

\end{columnas}
\end{document}
